\scanpage{28}
la condition de fidélité s'écrit
\[
\forall\ \text{diviseur }d\mid n,\quad\text{on a }\inserted{u_E^d\ne\mathrm{id}\ \text{i.e.}}\quad u_E^d(\deleted{\unsure{=}}s_0)\ne s_0
\tag*{(15)}
\]
(cf prop 4 (b)) -- et il suffit \inserted{\text{$=$}} de l'exiger pour $3\leq d\leq\left[\frac n3\right]$ (donc si $3\leq n\leq8$, cette condition est automatiquement satisfaite, le premier entier où elle ne l'est pas est $n=9$, où il suffit de \inserted{\text{\unsure{le}}} tester pour $n=9$ -- le $9$-gone régulier peut \deleted{\unsure{consis}} dégénérer pour consister trois fois un $3$-gône régulier --) -- Mais de $=$ qu'on peut laisser tomber cette condition de fidélité, on peut laisser tomber la \inserted{\text{dernière}} condition (14)
\[
(\sigma'_E\sigma_E)^{n}=\mathrm{id}_E
\]
ce qui signifie qu'on regarde en fait une représentation
\[
\varphi:D_\infty\mathbin{\unsure{-|\to}}\mathrm{Aff}(E)\tag*{(16)}
\]
du groupe diédral infini, sans se préoccuper si elle se factorise par $D_n$ ou non -- étant entendu que l'
% archive sheet number: 27
